\section*{Capítulo \ref{ch:b3:atmospheric-dynamics} --- Dinámica atmosférica}

\begin{solution}{exo:b3:atmospheric-dynamics:1}
(a)~$\zeta = 2\omega = 2\times2\pi/86\,400 = \qty{1.45e-4}{s^{-1}}$.
(b)~$f = \num{7.29e-5}$ y \qty{1.26e-4}{s^{-1}}: el vórtice tiene aproximadamente
el doble de la \omterm{def:b3:atmospheric-dynamics:vorticity-kinds}{vorticidad planetaria} a \ang{30} N y algo más que la de
\ang{60} N.
\end{solution}

\begin{solution}{exo:b3:atmospheric-dynamics:2}
$\zeta = f_{30} - f_{40} = \num{7.29e-5} - \num{9.37e-5} =
\qty{-2.08e-5}{s^{-1}}$: anticiclónico.
\end{solution}

\begin{solution}{exo:b3:atmospheric-dynamics:3}
$k = 2\pi/(5\times10^6) = \qty{1.257e-6}{m^{-1}}$; $c = -\beta/k^2 =
-\num{1.62e-11}/\num{1.58e-12} = \qty{-10.3}{m/s}$: hacia el oeste.
\end{solution}

\begin{solution}{exo:b3:atmospheric-dynamics:4}
$\beta = 2\Omega\cos\ang{60}/a = \qty{1.14e-11}{m^{-1}.s^{-1}}$;
$\lambda_s = 2\pi\sqrt{10/\num{1.14e-11}} = \qty{5870}{km}$; el círculo,
\qty{20015}{km}, admite $n_s = 3.4$ ondas.
\end{solution}

\begin{solution}{exo:b3:atmospheric-dynamics:5}
$q = f/h$ se conserva: sobre la cordillera $\zeta = f(8/10) - f =
\qty{-2.06e-5}{s^{-1}}$, anticiclónico; sobre la cuenca $f(12/10) - f =
\qty{+2.06e-5}{s^{-1}}$, ciclónico. Donde las columnas se estiran, a sotavento
de la cordillera y sobre la cuenca se genera vorticidad ciclónica: allí
las depresiones tienden a formarse.
\end{solution}

\begin{solution}{exo:b3:atmospheric-dynamics:6}
$k = 5/(a\cos\ang{45}) = \qty{1.11e-6}{m^{-1}}$, $\beta/k^2 =
\qty{13.1}{m/s}$: $c = \qty{1.9}{m/s}$, $c_g = \qty{28.1}{m/s}$. La energía
recorre \qty{5000}{km} en \num{2.1}~días, el patrón en \num{31}~días.
\end{solution}

\begin{solution}{exo:b3:atmospheric-dynamics:7}
(a)~$k^2 = \beta/U - l^2 = \num{1.079e-12} - \num{2.74e-13} =
\qty{8.05e-13}{m^{-2}}$, $k = \qty{8.97e-7}{m^{-1}}$: $\lambda =
\qty{7000}{km}$ y $n = 4.0$, menos ondas que con $l = 0$. (b)~$l^2 =
\num{1.097e-12} > \beta/U = \num{1.079e-12}$: $k^2$ sería negativo.
\end{solution}

\begin{solution}{exo:b3:atmospheric-dynamics:8}
$L_d = 0.012\times10^4/\num{1.031e-4} = \qty{1164}{km}$; $\sigma =
0.31\times\num{1.031e-4}\times\num{2e-3}/0.012 = \qty{5.3e-6}{s^{-1}}$, un
tiempo de $e$-folding de \num{2.2}~días; longitud de onda $2\pi\times1164/1.61 =
\qty{4540}{km}$.
\end{solution}

\begin{solution}{exo:b3:atmospheric-dynamics:9}
$[v^*T^*] = (VT_0/2)\cos\delta$: \qty{25}{K.m/s}, \qty{12.5}{K.m/s} y 0;
flujos de calor \num{25100}, \num{12600} y \qty{0}{W/m^2}. Cuando la
temperatura está desfasada un cuarto de longitud de onda respecto al viento, tanto
aire cálido va hacia el sur como hacia el norte: no hay transporte.
\end{solution}

\begin{solution}{exo:b3:atmospheric-dynamics:10}
Linealizando respecto a $U$ con $\zeta' = \nabla^2\psi'$ y $v' =
\partial_x\psi'$: $(\partial_t + U\partial_x)\nabla^2\psi' + \beta\partial_x
\psi' = 0$. Con $\psi' \propto \mathrm e^{\mathrm i(kx + ly - \omega t)}$,
$(\omega - Uk)K^2 = -\beta k$, así $\omega = Uk - \beta k/K^2$. Respecto
al aire, la velocidad de fase es $c - U = -\beta/K^2 < 0$ para toda onda.
\end{solution}

\begin{solution}{exo:b3:atmospheric-dynamics:11}
$\varphi_H = [5\times9.81\times12\,000\times0.25/(3\times(\num{7.292e-5})^2
\times(\num{6.371e6})^2)]^{1/2} = \qty{0.477}{rad}$, \ang{27.3}.
$u_M = \Omega a\varphi_H^2 = \qty{106}{m/s}$; exactamente, $\Omega a\sin^2
\varphi_H/\cos\varphi_H = \qty{110}{m/s}$: dos veces y media la
corriente en chorro observada.
\end{solution}

\begin{solution}{exo:b3:atmospheric-dynamics:12}
(a)~A través de la banda la interfase se desplaza una cantidad $\eta = sy$, $-L/2 \le y
\le L/2$, desde su posición horizontal. Una capa con contraste de densidad reducido
desplazada una cantidad $\eta$ almacena $\tfrac12\rho g'\eta^2$ por unidad de área; promediada
sobre la banda, $\langle\eta^2\rangle = s^2L^2/12$. (b)~$g' = 0.327$:
\num{1.5e5} y \qty{6.1e5}{J/m^2}, frente a $\tfrac12\times0.7\times15^2
\times10^4 = \qty{7.9e5}{J/m^2}$ de energía cinética: comparables.
\end{solution}

\begin{solution}{pb:b3:atmospheric-dynamics:1}
\textbf{1.} $f_{45} = \qty{1.031e-4}{s^{-1}}$, $f_{50} =
\qty{1.117e-4}{s^{-1}}$; $\zeta = f_{45} - f_{50} = \qty{-8.6e-6}{s^{-1}}$.

\textbf{2.} $R = V/|\zeta| = 15/\num{8.6e-6} = \qty{1750}{km}$; en sentido horario,
de regreso hacia el sur.

\textbf{3.} $\zeta = f_{45}(9.5/10) - f_{50} = \qty{-1.38e-5}{s^{-1}}$, más
anticiclónico.

\textbf{4.} $\zeta = f_{45} - f_{40} = \qty{+9.4e-6}{s^{-1}}$, ciclónico: la parcela
gira hacia el norte. Desplazada en uno u otro sentido, la parcela adquiere la vorticidad
que la hace regresar: un efecto restaurador, y una oscilación.

\textbf{5.} $\beta = \qty{1.62e-11}{m^{-1}.s^{-1}}$; el círculo mide
\qty{28306}{km} de longitud.

\textbf{6.} $n = 3$ a 8: longitudes de onda 9435, 7076, 5661, 4718, 4044 y
\qty{3538}{km}; $c = -21.5$, $-5.5$, 1.9, 5.9, 8.3 y \qty{9.9}{m/s}; $c_g
= 51.5$, 35.5, 28.1, 24.1, 21.7 y \qty{20.1}{m/s}.

\textbf{7.} $n = 3$ y 4 se desplazan hacia el O, $n \ge 5$ hacia el E. La onda $n = 5$ tarda
$10^6/1.86 = \qty{5.4e5}{s}$, \num{6.2}~días, en recorrer \qty{1000}{km}.

\textbf{8.} Aguas abajo, hacia el E: a $c_g = \qty{24.1}{m/s}$ la energía recorre
\qty{5000}{km} en \num{2.4}~días.

\textbf{9.} $\lambda_s = 2\pi\sqrt{15/\num{1.62e-11}} = \qty{6050}{km}$,
$n_s = 4.7$.

\textbf{10.} $n_s = 5.7$ para \qty{10}{m/s} y 4.1 para \qty{20}{m/s}: un
viento del oeste más intenso favorece ondas estacionarias más largas, con menos
que dan la vuelta al hemisferio.

\textbf{11.} $n_s = 5.8$ a \ang{35} N y 3.4 a \ang{55} N.

\textbf{12.} $\zeta = f(8/10) - f = \qty{-2.06e-5}{s^{-1}}$: la compresión sobre
las montañas hace que la columna sea anticiclónica; el estiramiento a sotavento
la vuelve ciclónica. Las montañas imponen al flujo un patrón fijo de vorticidad,
un forzamiento que no se desplaza.

\textbf{13.} Un forzamiento fijo produce una respuesta permanente; la respuesta es
máxima cuando el patrón forzado es en sí mismo una onda libre que estaría
quieta, la resonancia en $\lambda_s$.

\textbf{14.} $\Lambda = 25/10^4 = \qty{2.5e-3}{s^{-1}}$;
$|\partial T/\partial y| = \Lambda fT/g = \num{2.5e-3}\times\num{1.031e-4}
\times250/9.81 = \qty{6.6}{K}$ por \qty{1000}{km}.

\textbf{15.} $\mathrm{Ri} = 0.011^2/(\num{2.5e-3})^2 = 19$: muy por encima de 0.25,
sin turbulencia por cizalladura.

\textbf{16.} $L_d = 0.011\times10^4/\num{1.031e-4} = \qty{1067}{km}$.

\textbf{17.} $2\pi\times1067/1.61 = \qty{4160}{km}$; $28306/4160 = 6.8$
ondas.

\textbf{18.} $2\pi\times1067/2.40 = \qty{2790}{km}$.

\textbf{19.} $\sigma = 0.31\times\num{1.031e-4}\times\num{2.5e-3}/0.011 =
\qty{7.3e-6}{s^{-1}}$, \qty{0.63}{\text{día}^{-1}}.

\textbf{20.} $\exp(3\times0.627) = 6.6$.

\textbf{21.} La tormenta, \qty{4160}{km}, es más corta que la \omterm{prop:b3:atmospheric-dynamics:stationary}{longitud de onda
estacionaria de Rossby}, \qty{6050}{km}: como onda de Rossby, la arrastra hacia el este el
viento del oeste, y las tormentas avanzan hacia el este a través de las ondas largas
más lentas o estacionarias.

\textbf{22.} Con la altura, hacia el oeste: solo entonces el aire que se desplaza hacia el polo es
más cálido que el aire que se desplaza hacia el ecuador, de modo que el calor se dirige hacia el polo y la
\omterm{def:b3:atmospheric-dynamics:ape}{energía potencial disponible} se libera.

\textbf{23.} La perturbación de crecimiento más rápido se amplifica en un factor $e$ cada
\textbf{\num[text-series-to-math=true]{1.6}~días} (\num{1.59}~días según el
cálculo).
\end{solution}
